\section{Verification purpose}

This case verifies coupling among the simulation clock, solar burn-period calculation, and surface-fire diurnal multiplier. With uniform fuel, moisture, wind, and slope, the question is whether ELMFIRE reduces and restores radial advance at the prescribed clock hours, each within \SI{1}{hour}. The case exposes clock offsets, failed midnight wrapping, reversed interval tests, or an unapplied multiplier; it does not calibrate the multiplier.

The scientific setup and historical expected response are reconstructed from Section~3.2 of the supplied \emph{ELMFIRE Guide}. The Guide is used as reference provenance, not as an instruction or a substitute for evidence. Configuration statements below are checked against the generated GeoTIFFs and namelist; current results and decisions come only from case-specific outputs and generated metrics. Where the generated experiment differs from the historical example, the report states that difference explicitly.

\section{Mathematical formulation and numerical implementation}

For simulation time $t$ seconds and forecast start $H_0$, ELMFIRE uses
\begin{equation}H(t)=\operatorname{mod}\!\left(H_0+t/3600,24\right).\end{equation}
The solar routine calculates sunrise and sunset from the lower-left coordinate, year, and hour of year. With default center fraction $f_c=0.667$ and burn length $L_b=\SI{10}{hour}$,
\begin{align}H_c&=H_{\rm rise}+f_c(H_{\rm set}-H_{\rm rise}),\\H_b&=H_c-L_b/2,&H_e&=H_c+L_b/2.\end{align}
The stop hour is wrapped above 24. For this non-wrapping interval,
\begin{equation}D(t)=\begin{cases}1,&H_b\le H(t)\le H_e,\\0.1,&H(t)<H_b\ \text{or}\ H(t)>H_e.\end{cases}\end{equation}
The spread-rate calculation multiplies its unadjusted rate by $D(t)$. Uniform flat, wind-free conditions produce an approximately circular fire whose equivalent radius $r(t)=\sqrt{A(t)/\pi}$ has full and reduced slope plateaus.

\section{Derivation of expected behavior}

The controlled response specifies burn-period start \SI{7.5}{hour} and stop \SI{17.5}{hour}. Full-to-reduced and reduced-to-full transitions are therefore expected at those respective clock hours. Since $H_0=\SI{10}{hour}$, the events first occur approximately \SI{7.5}{hour} and \SI{21.5}{hour} after ignition and repeat during the three-day run. This reference calculation follows from the controlled targets and clock equations; calculated times are independently inferred from the arrival raster.

% BEGIN EXPLICIT METRIC DERIVATION
For the generated astronomical inputs, local sunrise and sunset are
approximately \(H_r=4.5\) h and \(H_s=16.5\) h.  The operational cutoff
with fractional daylight \(f_c=0.667\) is
\[
 H_c=H_r+f_c(H_s-H_r)=4.5+0.667(16.5-4.5)
 =12.504~\mathrm{h}\simeq12.5~\mathrm{h}.
\]
A ten-hour operating window centered on \(H_c\) therefore has bounds
\(H_c\pm5=\{7.5,17.5\}\) h.  Relative to the 10:00 forecast start, the
first reduction occurs after
\((17.5-10)(3600)=27000~\mathrm{s}\); after midnight wrapping, full
operation resumes after \((24-10+7.5)(3600)=77400~\mathrm{s}\).
The transition metrics compare detected changes with these reference
times using the tabulated timing tolerance.
% END EXPLICIT METRIC DERIVATION

\section{Verification metrics and acceptance criteria}

The evaluator selects the lexically latest fire arrival time GeoTIFF, masks nodata and values outside $0\le t<10^{8}$ seconds, and requires 50 burned cells. Analysis ends at the earliest domain-edge arrival so boundary truncation cannot mimic slowdown. At \SI{30}{minute} intervals, burned-cell count is converted to equivalent radius and differenced. After radius exceeds three cells, rates below half the 80th percentile of positive rates are reduced. Median full-to-reduced and reduced-to-full clock hours are then calculated.

\begin{center}\begin{tabular}{@{}llll@{}}
\toprule Metric & Target & Acceptance rule & Role \\
\midrule
Reduction starts & \SI{17.5}{hour} & absolute error $\le\SI{1}{hour}$ & burn-period stop \\
Reduction ends & \SI{7.5}{hour} & absolute error $\le\SI{1}{hour}$ & burn-period start \\
\bottomrule\end{tabular}\end{center}

Both rows are required. A missing raster, insufficient burned area or duration, or undetected transition produces NOT EVALUABLE, never PASS.

\section{Simulation configuration and predicted outcome}

\begin{center}\begin{tabular}{@{}lll@{}}
\toprule Input & Value & Verification role \\
\midrule
Domain/grid & $126\times126$; \SI{30}{m} & circular-front resolution \\
CRS & EPSG:32635 & location for solar calculation \\
Fuel/terrain & uniform FM3; zero slope & constant baseline spread \\
Wind/dead moisture & zero; 6, 7, 8 percent & constant forcing \\
Live moisture & herb 60; woody 90 percent & constant forcing \\
Ignition & $(1890,1890)$ m at $t=0$ & centered radial growth \\
Time control & \SI{1}{s}; \SI{259200}{s} stop & three daily cycles \\
Calendar & year 2025; hour 2890; start 10 & clock coupling \\
Diurnal control & enabled; factor 0.1 & two rate plateaus \\
Output & final time of arrival (TOA), spread rate, FLIN & case-specific evidence \\
\bottomrule\end{tabular}\end{center}

The predicted arrival field is concentric; contour spacing contracts overnight and expands during the burn period, with both transitions inside their acceptance bands.

Figure~\ref{fig:input-configuration} shows the prescribed spatial condition used in this experiment. The displayed field represents the initial level-set geometry or a controlling fuel, terrain, moisture, wind, or canopy quantity, as identified by the figure labels. The plotted values come from the prepared simulation inputs, not from the expected response.

\begin{figure}[H]
  \centering
  \IfFileExists{../figures/input_configuration.pdf}{%
    \includegraphics[width=0.98\textwidth,height=0.46\textheight,keepaspectratio]{input_configuration.pdf}%
  }{\fbox{\parbox{0.88\textwidth}{Prepared-input evidence is unavailable. The prescribed configuration is described in the text; this absence does not establish that it was simulated.}}}
  \caption{Prepared input conditions from the prescribed spatial fields. Coordinates are in metres; the colorbar identifies the displayed field and its physical units or category. The ignition panel is derived directly from \(\phi\leq0\).}
  \label{fig:input-configuration}
\end{figure}

\section{Actual simulation results}

The current calculated status is \textbf{\VerificationStatus};
\PassingMetricCount{} of \MetricCount{} required comparisons pass across
\VariantCount{} completed simulation conditions. The displayed measurements
come from these simulations, not from the historical Guide. Missing evidence
would remain NOT EVALUABLE.

\begin{figure}[H]
\centering
\IfFileExists{../figures/whole_domain_result.pdf}{\includegraphics[width=0.82\textwidth,height=0.58\textheight,keepaspectratio]{whole_domain_result.pdf}}{\fbox{\parbox{0.85\textwidth}{No case-specific whole-domain arrival-time raster is available.}}}
\caption{Final whole-domain time of arrival, in seconds, when a case-specific result figure has been generated. The figure must identify its selected GeoTIFF and simulation condition; the generated metrics table, not visual appearance, determines acceptance.}
\end{figure}

\section{Calculated metrics and verification decision}

\begin{center}\begin{tabular}{@{}llll@{}}
\toprule Simulation condition / metric & Expected & Calculated & Status \\
\midrule
\MetricRows
\bottomrule\end{tabular}\end{center}

The table and the decision \textbf{\VerificationStatus} report the existing
calculated metrics. Missing required evidence yields NOT EVALUABLE; no
metric values are inferred from the appearance of the simulated fire.

\section{Technical interpretation}

The recorded transition measurements satisfy the stated one-hour bands,
supporting the coupling between simulation time, the solar window, and the
diurnal spread factor for the prescribed conditions. A common shift in both
transitions would instead suggest an error in the clock, calendar, coordinates,
or solar timing; missing contrast would suggest an unapplied multiplier or
insufficient simulation history before the transition. This interpretation
uses the established calculated metrics and does not extend the claim beyond
the configured experiment.

\section{Scope and limitations}

This synthetic case covers one date, location, fuel, factor, and grid. It does not validate factor magnitude, other seasons or coordinate systems, a burn interval wrapping midnight, heterogeneity, wind, or crown-fire, spotting, suppression, and interpolation interactions. The estimator assumes circular growth, discards post-edge behavior, samples every \SI{30}{minute}, and infers rather than records the internal factor. This is implementation verification, not validation of real-fire nocturnal behavior.
